%% GENERATED by tools/from_docs.py from stepss-docs. Do not edit by hand:
%% edit the page named below in stepss-docs and regenerate.
%% source: stepss-docs/src/content/docs/python/helios.md

\chapter{Power flow from Python}\label{doc:py-helios}

The \texttt{stepss.helios} module wraps \hyperref[ch-pf]{Helios}, the STEPSS AC power-flow engine. Pre-compiled libraries are bundled with the package for Windows, Linux, and macOS, so no separate installation is required.

Unlike the RAMSES classes (which follow the historical camelCase conventions), this module uses PEP 8 snake\_case naming. Errors are raised as \texttt{stepss.HeliosError}, carrying the engine's diagnostic message.

\section{Basic Workflow}\label{py-helios:basic-workflow}

\texttt{HeliosSession} follows a load \ensuremath{\rightarrow} (modify) \ensuremath{\rightarrow} solve \ensuremath{\rightarrow} query lifecycle and works as a context manager:

\begin{Verbatim}
from stepss.helios import HeliosSession

with HeliosSession() as pf:
    pf.load_file('network.dat')
    converged = pf.solve()

    v, angle = pf.get_bus_voltage('1041')     # one bus (pu, rad)
    p, q = pf.get_branch_flow('1042-1044')    # from-end flow (MW, Mvar)
\end{Verbatim}

Solve diagnostics are available after \texttt{solve()}: \texttt{pf.converged}, \texttt{pf.solver\_status}, \texttt{pf.iterations}, and \texttt{pf.max\_mismatch}.

\section{Solver Options}\label{py-helios:solver-options}

Options map to the \texttt{\$PARAM} records of the data file and can be overridden after loading (loading overwrites them with the file's values, so set options \emph{after} \texttt{load\_file}):

\begin{Verbatim}
from stepss.helios import Option

pf.load_file('network.dat')
pf.set_option(Option.TOLAC, 0.001)     # MW
pf.set_option(Option.MAX_ITER, 30)
pf.solve()
\end{Verbatim}

\section{Vectorized Results}\label{py-helios:vectorized-results}

Bulk getters return NumPy arrays indexed like the corresponding \texttt{*\_names()} lists, with no text parsing needed:

\begin{Verbatim}
v_pu, angle_rad = pf.get_bus_voltages()
p_load, q_load = pf.get_bus_loads()
p_from, q_from, p_to, q_to = pf.get_branch_flows()
p_gen, q_gen, status = pf.get_generator_outputs()

names = pf.bus_names()
print(f"lowest voltage: {v_pu.min():.4f} pu at {names[v_pu.argmin()]}")
\end{Verbatim}

Per-element dataclasses are available via \texttt{get\_bus\_info()}, \texttt{get\_branch\_info()}, \texttt{get\_generator\_info()}, and \texttt{get\_converter\_info()}.

\section{AC/DC Converters}\label{py-helios:acdc-converters}

A case with \texttt{BUSDC} and \texttt{ACDC} records (see \hyperref[ch-pf]{AC/DC Converters (ACDC)}) solves like any other, and its DC buses appear in \texttt{bus\_names()} with \texttt{bus\_type} \texttt{BusType.DC}. The converters have their own getters, which need \texttt{stepss} 3.85.2 or later (Helios 1.7.0):

\begin{Verbatim}
from stepss.helios import ConverterMode

with HeliosSession() as pf:
    pf.load_file('case5_mtdc_droop.dat')
    pf.solve()

    for name in pf.converter_names():
        c = pf.get_converter_info(name)
        print(f"{c.name} {c.mode.name:5} PAC {c.pac_mw:8.2f} MW  "
              f"PDC {c.pdc_mw:8.2f} MW  loss {c.converter_loss_mw:5.2f} MW")

    pac, qac, pdc, status = pf.get_converter_outputs()   # NumPy arrays
\end{Verbatim}

\texttt{get\_converter\_info()} returns a \texttt{ConverterInfo} with the same quantities as the converter table of the power-flow menu:

\begin{small}\begin{longtable}[]{@{}
  >{\raggedright\arraybackslash}p{(\linewidth - 4\tabcolsep) * \real{0.3465}}
  >{\raggedright\arraybackslash}p{(\linewidth - 4\tabcolsep) * \real{0.6535}}
@{}}
\toprule\noalign{}
\begin{minipage}[b]{\linewidth}\raggedright
Field
\end{minipage} & \begin{minipage}[b]{\linewidth}\raggedright
Meaning
\end{minipage} \\
\midrule\noalign{}
\endhead
\bottomrule\noalign{}
\endlastfoot
\texttt{ac\_bus}, \texttt{dc\_bus}, \texttt{in\_service} & Connection and status \\
\texttt{pac\_mw}, \texttt{qac\_mvar} & Power at the AC bus: \texttt{PAC} drawn from it, \texttt{QAC} injected into it \\
\texttt{pdc\_mw} & Power delivered into the DC bus; \texttt{PAC} and \texttt{PDC} are positive when rectifying \\
\texttt{mode} & \texttt{ConverterMode}: \texttt{PQ}, \texttt{PV}, \texttt{VDC}, \texttt{DROOP}, \texttt{LIMIT} or \texttt{REF}, as the table's mode column \\
\texttt{q\_control} & \texttt{-2} holding \texttt{VAC}, \texttt{1} fixed Q, \texttt{2} at \texttt{QMAX}, \texttt{3} at \texttt{QMIN} \\
\texttt{rectifying} & Direction of the active power \\
\texttt{vc\_pu}, \texttt{pc\_mw}, \texttt{qc\_mvar} & Converter node behind the station; without a station they equal the AC bus voltage, \texttt{PAC} and \texttt{-QAC} \\
\texttt{station\_loss\_mw}, \texttt{converter\_loss\_mw} & \texttt{PAC\ -\ Pc} lost in the station and \texttt{Pc\ -\ PDC} in the converter \\
\end{longtable}\end{small}

An out-of-service converter reports zero powers. A case that fails on its DC side ends with \texttt{pf.solver\_status} set to \texttt{SolverStatus.NO\_DC\_REGULATION} or \texttt{SolverStatus.CONVERTER\_INFEASIBLE}, and \texttt{pf.last\_error} names the element (see \hyperref[ch-pf]{Solver Outcomes}).

With an older bundled Helios, \texttt{pf.has\_converter\_api} is \texttt{False} and the converter getters raise \texttt{HeliosError}; the rest of the session works as before.

\section{Modifying the System}\label{py-helios:modifying-the-system}

Modifications accumulate an active-power imbalance that \texttt{apply\_changes()} settles: connectivity check, redispatch onto the remaining generators, and a re-solve (the same workflow as the interactive modify menu). \texttt{change\_*} methods apply increments; \texttt{set\_load()} sets absolute values:

\begin{Verbatim}
pf.trip_branch('1042-1044')
pf.change_load('1041', 50.0, 10.0)     # +50 MW, +10 Mvar
pf.set_generator_voltage('g6', 1.02)
pf.apply_changes()                     # redispatch + re-solve

pf.reset()                             # back to the state saved at load
\end{Verbatim}

\section{Contingency Screening}\label{py-helios:contingency-screening}

Run automatic N-1 over selected equipment classes, or a Fortran-format contingency file (\texttt{BT}/\texttt{GT}/\texttt{ST}/\texttt{HC} actions). Each result is a dataclass with convergence, acceptance, voltage/loading extremes, and violation strings; the base case is untouched afterwards:

\begin{Verbatim}
results = pf.run_contingencies(branches=True, generators=True,
                               v_min=0.95, v_max=1.05, overload=100.0)
for r in results:
    if not r.accepted:
        print(r.name, r.violations)

results = pf.run_contingencies(file='contingencies.txt')
\end{Verbatim}

\section{Exporting Results}\label{py-helios:exporting-results}

\begin{Verbatim}
pf.write_dump('solved_case.dat')      # re-loadable data file
pf.write_voltrat('volt_rat.dat')      # LFRESV + TRANSFO records, RAMSES initial conditions
pf.write_matlab('system.m')           # operating point + Y-bus script
pf.write_diagram('template.svg', 'diagram.svg')
\end{Verbatim}

\texttt{volt\_rat.dat} is the natural bridge to dynamic simulation: solve a power flow with Helios, export it, and pass it to a RAMSES case via \texttt{case.addData(\textquotesingle{}volt\_rat.dat\textquotesingle{})}.

\begin{docsnote}{Note}

\texttt{write\_voltrat()} is the API form of the \texttt{VT} menu command, and writes the same
records. For what those records contain, and how they relate to the hand-written
\texttt{TRFO} style, see
\hyperref[ch-pf]{Exported Operating Point}.

\end{docsnote}

\section{Runnable Examples}\label{py-helios:runnable-examples}

Five self-contained scripts ship in the repository under \href{https://github.com/SPS-L/stepss-python-ui/tree/master/examples/helios}{\texttt{examples/helios/}}: basic solve, options and arrays, modify and re-solve, contingency screening, and exports.

\section{Further Reading}\label{py-helios:further-reading}

\begin{itemize}
\tightlist
\item
  \hyperref[ch-pf]{Power Flow user guide}, data format, solver parameters, engine details
\item
  Full method-level documentation: the docstrings on \texttt{HeliosSession}, e.g.~\texttt{help(stepss.HeliosSession)}
\end{itemize}
